%% kapa: exam=1 round=20 year=2009 period=2 form=A subject=appraisal_law subject_ko=부동산관계법규 seq=18 number=58
%% source: Q-Net 감정평가사 기출문제 (cst00309) artlSeq=5024178 file=extracted/1008/5024178/2009년 제20회 감정평가사 1차시험 기출문제/2교시.hwp
%% conversion: kordoc 4.24.51 -> tools/gen1.py (auto); review=proofread-sonnet; flags=none
\begin{question}{58}
부동산 가격 공시 및 감정평가에 관한 법령상 개별공시지가의 결정․공시에 관한 설명으로 옳지 \uline{않은} 것은?
\begin{choices}
  \item 국토해양부장관은 매년 공시지가의 공시기준일 현재 관할구역 안의 개별토지의 단위면적당 가격을 결정․공시하고, 이를 관계행정기관 등에 제공하여야 한다.
  \item 조세 또는 부담금 등의 부과대상이 아닌 토지인 경우에는 개별공시지가를 결정․공시하지 아니할 수 있다.
  \item 개별공시지가에 대하여 이의가 있는 자는 개별공시지가의 결정ㆍ공시일부터 30일 이내에 서면으로 이의를 신청할 수 있다.
  \item 개별공시지가를 결정․공시하는 경우에는 당해 토지와 유사한 이용가치를 지닌다고 인정되는 하나 또는 둘 이상의 표준지의 공시지가를 기준으로 토지가격비준표를 사용하여 지가를 산정하여야 한다.
  \item 개별공시지가를 결정․공시하기 위하여 개별토지의 가격을 산정한 때에는 원칙적으로 그 타당성에 대하여 감정평가업자의 검증을 받고 토지소유자 그 밖의 이해관계인의 의견을 들어야 한다.
\end{choices}
\end{question}
\answer{1}
